% Figure: the four bolted-joint failure modes of a single-bolt lap joint,
% drawn as four small plan/section sketches.
\begin{figure}[htbp]
\centering
\begin{tikzpicture}[>=latex, scale=1.0]
  % ==== 1: bolt shear ====
  \begin{scope}
    \fill[engblue!15, draw=engblue, thick] (-1.3,0.15) rectangle (1.3,0.55);   % top plate
    \fill[engblue!25, draw=engblue, thick] (-1.3,-0.55) rectangle (1.3,-0.15);  % bottom plate
    \fill[enggray] (-0.12,-0.55) rectangle (0.12,0.55);                          % bolt
    \draw[->, thick, engred] (-1.3,0.35) -- (-1.8,0.35);
    \draw[->, thick, engred] (1.3,-0.35) -- (1.8,-0.35);
    % shear plane
    \draw[engred, very thick] (-0.18,0) -- (0.18,0);
    \node[font=\scriptsize, anchor=north] at (0,-0.7) {bolt shear};
  \end{scope}
  % ==== 2: bearing ====
  \begin{scope}[xshift=4cm]
    \fill[engblue!18, draw=engblue, thick] (-1.3,-0.55) rectangle (1.3,0.55);
    % elongated hole
    \fill[white, draw=engblue] (-0.35,-0.3) rectangle (0.35,0.3);
    \fill[enggray] (0.0,0) circle (0.28);
    \draw[->, thick, engred] (0.0,0) -- (0.75,0);
    \node[engred, font=\scriptsize, anchor=west] at (0.5,0.28) {bears};
    \node[font=\scriptsize, anchor=north] at (0,-0.7) {bearing};
  \end{scope}
  % ==== 3: net-section tension ====
  \begin{scope}[xshift=8cm]
    \fill[engblue!18, draw=engblue, thick] (-1.3,-0.55) rectangle (1.3,0.55);
    \fill[white, draw=engblue] (0,0) circle (0.28);
    % tear lines through net section
    \draw[engred, very thick] (0,0.28) -- (0,0.55);
    \draw[engred, very thick] (0,-0.28) -- (0,-0.55);
    \draw[->, thick, engred] (-1.3,0) -- (-1.8,0);
    \draw[->, thick, engred] (1.3,0) -- (1.8,0);
    \node[font=\scriptsize, anchor=north] at (0,-0.7) {net section};
  \end{scope}
  % ==== 4: block shear ====
  \begin{scope}[xshift=12cm]
    \fill[engblue!18, draw=engblue, thick] (-1.3,-0.55) rectangle (1.3,0.55);
    \fill[white, draw=engblue] (0.3,0) circle (0.25);
    % the block that tears out (behind bolt toward right edge)
    \draw[engred, very thick] (0.3,0.25) -- (1.3,0.25);   % shear plane top
    \draw[engred, very thick] (0.3,-0.25) -- (1.3,-0.25);  % shear plane bottom
    \draw[engred, very thick] (0.05,0.25) -- (0.05,-0.25); % tension plane
    \draw[->, thick, engred] (-1.3,0) -- (-1.8,0);
    \node[font=\scriptsize, anchor=north] at (0,-0.7) {block shear};
  \end{scope}
\end{tikzpicture}
\caption[Four bolted-joint failure modes]{The four ways a single-bolt
lap joint in tension fails, each invisible to a global free-body diagram
that treats the connection as a frictionless pin. Bolt shear: the bolt
parts across the plane between the plates. Bearing: the bolt shank crushes
and elongates the hole until the plate ahead of it tears. Net-section
tension: the plate, weakened by the hole, fractures across the reduced
area between hole and edge. Block shear: a chunk of plate behind the bolt
tears out along two shear planes parallel to the load and one tension
plane across it, the mechanism the NBS investigation named for the
modified upper connection in the Hyatt Regency walkway
\cite{acc:nbs-hyatt-1982}. A connection check is incomplete until all four
capacities are computed and the smallest is taken as the joint's
strength.}
\label{fig:vol03ch02:bolt-failure-modes}
\end{figure}
